\subsection*{B.2 Entrevista N.°~2 --- Morena Cejas}
\addcontentsline{toc}{subsection}{B.2 Entrevista N.° 2 --- Morena Cejas}

\noindent\textbf{Entrevistada:} Morena Cejas\\
\noindent\textbf{Perfil:} Dueña de emprendimiento de indumentaria (ropa unisex: remeras, bermudas, buzos, joggings). Vende por Instagram, boca en boca y Tienda Nube. Tiene menos de un año de antigüedad.\\
\noindent\textbf{Entrevistadores:} Tomás Gonçalves (moderador) y Micaela Palomino (observadora)\\
\noindent\textbf{Fecha:} Sábado 4 de julio de 2026\\
\noindent\textbf{Modalidad:} Videollamada (semiestructurada)\\
\noindent\textbf{Plataforma del negocio:} Tienda Nube (plan gratuito al momento de la entrevista; previamente pagó el primer nivel)

\begin{framed}
\small\noindent\textit{Nota metodológica: Esta transcripción fue limpiada a partir de la grabación original. Se eliminó charla trivial previa y posterior a la entrevista, se corrigió la sintaxis para mejorar la fluidez de lectura y se reasignaron las atribuciones de orador donde la herramienta de transcripción automática las asignó erróneamente (verificado por contexto). No se agregó, inventó ni dedujo información que los participantes no hayan dicho explícitamente.}
\end{framed}

\subsubsection*{Bloque 0 · Introducción y rapport}

\noindent\textbf{Tomás Gonçalves:} Estuve viendo la página que tenés. Nos viene re bien porque está hecha con Tienda Nube. Te quería preguntar: veo que vendés ropa nada más. ¿Hace cuánto arrancaste más o menos con esto?

\noindent\textbf{Morena Cejas:} Hace casi un año. Todavía no llegué al año; me faltan tres meses para cumplirlo.

\subsubsection*{Bloque 1 · Contexto del negocio y herramientas actuales}

\noindent\textit{--- Registro de ventas y uso de Tienda Nube ---}

\noindent\textbf{Micaela Palomino:} Entiendo que hace un año tenés este emprendimiento de ropa. ¿Dónde registrás todas las ventas?

\noindent\textbf{Morena Cejas:} Tengo armado un Excel y anoto todo ahí. Igual, en Tienda Nube también se registran. Pero hasta ahora Tienda Nube fue lo que menos usé, porque en general les vendí a conocidos o amigos de conocidos. Entonces, para que se ahorren ese paso o el envío, fue más hablar por chat y arreglar para llevárselos.

\noindent\textbf{Micaela Palomino:} ¿Tienda Nube te cobra comisión por venta?

\noindent\textbf{Morena Cejas:} Depende del método de pago. Lo mejor que podés poner es transferencia, que es la que menos te cobra. Yo elegí transferencia porque era lo que menos me sacaba de mi ganancia.

\noindent\textbf{Micaela Palomino:} Vos vas registrando todas las ventas en Excel. ¿Después qué hacés con el Excel? ¿Te sirve para ver el stock?

\noindent\textbf{Morena Cejas:} Me sirve para ver stock. También anoto cuánto me salió la prenda y a cuánto la vendí, porque hay prendas que vendí en promoción, alguna liquidación o algún descuento, y lo anoté ahí como para saber cuánto gané de esa prenda.

\noindent\textbf{Micaela Palomino:} O sea, todo el análisis y las métricas básicamente las sacás del Excel y las llevás vos.

\noindent\textbf{Morena Cejas:} Sí, ya te digo, la mayor parte de las ventas las hice por mi cuenta, entonces Tienda Nube no la utilicé tanto en ese sentido. Pero tengo entendido que igual te da un porcentaje de ventas, cuánta gente entró, cuánta gente se quedó en el carrito y no siguió la compra; eso te hace Tienda Nube como registro.

\noindent\textit{--- Métricas de Tienda Nube y carritos abandonados ---}

\noindent\textbf{Tomás Gonçalves:} ¿Qué es lo que más te sirve de esas métricas?

\noindent\textbf{Morena Cejas:} A mí me sirve saber cuánta gente se quedó en el carrito y no siguió la compra, porque además te da la opción de generar un mensaje: escribirle a la persona una vez que pone sus datos en Tienda Nube. Cuando te quedás en el carrito, te manda algo como ``podés seguir tu compra'' o ``por qué te quedaste ahí''.

\noindent\textbf{Micaela Palomino:} ¿Esa opción es paga o solamente la tenés que configurar?

\noindent\textbf{Morena Cejas:} Solo la tenés que configurar. Pero Tienda Nube tiene varios planes. Yo al principio empecé con el gratuito, y después, para incorporar distintos diseños para la página y métodos de pago y envío, pagué un plan más que no era tan caro. Te da la opción de Andreani y Correo Argentino para tener las dos, porque la gente a veces prefiere uno u otro, y lo mismo para las formas de pago.

\noindent\textbf{Micaela Palomino:} Cuando Tienda Nube manda la notificación a la persona sobre el carrito, ¿le dice lo que tiene específicamente?

\noindent\textbf{Morena Cejas:} Sí, te dice lo que tenés específicamente.

\noindent\textbf{Micaela Palomino:} ¿Y a vos te da una métrica que te diga, por ejemplo, que tal producto se quedó mucho en el carrito?

\noindent\textbf{Morena Cejas:} No, eso no. Cuando la gente se queda en el carrito, tampoco me avisaba. Era si me metía en la aplicación; ahí estaba la opción de ver.

\subsubsection*{Bloque 2 · Gestión de inventario y stock}

\noindent\textit{--- Notificaciones de stock y control manual ---}

\noindent\textbf{Tomás Gonçalves:} Nosotros nos reunimos con una emprendedora que también tiene un emprendimiento, y con el tema de stock se entera en el momento cuando se queda sin nada. ¿Vos más o menos tenés una proyección, o te enterás de golpe?

\noindent\textbf{Morena Cejas:} El tema es que Tienda Nube no te avisa. Sí se va registrando, pero no te manda una notificación. Yo entraba a la página y lo iba registrando, ponele, de las ventas que hacía yo por mi cuenta. Pero no es que tenés muchas notificaciones de Tienda Nube.

\noindent\textbf{Micaela Palomino:} ¿Te acordás alguna vez de quedarte sin stock de algo que lo estabas vendiendo bien?

\noindent\textbf{Morena Cejas:} A ver, como ya te digo, no hice tantas ventas por Tienda Nube. No me pasó quedarme sin stock en Tienda Nube porque me compraron por ahí, sino que me quedé sin stock por las ventas por mi cuenta. Además, como la tengo yo, la ropa está en mi casa, me voy dando cuenta.

\noindent\textbf{Micaela Palomino:} O sea, es más por control visual, porque la tenés en tu casa.

\noindent\textbf{Morena Cejas:} Sí, por las dudas, por el tema de los talles. Si me quedé sin algún talle, lo tengo para fijarme, pero no es que Tienda Nube me avisa.

\noindent\textit{--- Reposición de stock y productos no vendidos ---}

\noindent\textbf{Micaela Palomino:} Más allá de Tienda Nube, en tu operatoria en general, ¿alguna vez repusiste stock de algo que al final no vendiste?

\noindent\textbf{Morena Cejas:} Sí, de remeras. O de algún producto muy específico. Pero en general no repuse tanto stock de una cosa, sino que metí producto nuevo. Me pasó por un producto nada más.

\noindent\textbf{Micaela Palomino:} ¿Por qué pensás que no se vendió?

\noindent\textbf{Morena Cejas:} Porque eran las bermudas y se pasó la temporada del calor. Yo empecé en octubre-noviembre y ahí es donde más se vendían las bermudas. Después, ya no.

\noindent\textit{--- Restricciones de producción con fábricas ---}

\noindent\textbf{Micaela Palomino:} ¿Pero no podés hacer, por ejemplo, diez negros y dos rojos?

\noindent\textbf{Morena Cejas:} El tema es que, por lo menos con la fábrica con la que trabajo yo, primero tiene un mínimo de prendas que podés hacer ---no sé, doce de un modelo---. Yo puedo hacer doce buzos o doce joggings, pero los colores y las telas también varían: el negro lo puedo hacer tranquilamente, pero para el rojo tengo que pedir un rollo entero, entonces tengo que hacer todo rojo sí o sí. Y el resto de las fábricas también tienen otros mínimos; hay fábricas que piden mínimos de cien prendas. No me voy a tirar con cien prendas, es un montón.

\noindent\textbf{Micaela Palomino:} Es un montón, y sobre todo si no tenés publicidad. Siento que si tenés muchos seguidores en TikTok y metés un montón de marketing, un poco te sale.

\noindent\textbf{Morena Cejas:} Claro. Es mi idea, meterle al marketing y la publicidad, porque creo que me va a ayudar mucho más.

\subsubsection*{Bloque 3 · Análisis de ventas y decisiones comerciales}

\noindent\textit{--- Publicidad y canales de venta ---}

\noindent\textbf{Tomás Gonçalves:} ¿Hacés publicidad?

\noindent\textbf{Morena Cejas:} Hago contenido más que nada en redes para que se mueva. Pero en un momento me quedé medio estancada con el contenido porque me bajoné: no era lo que yo esperaba, y perdí mucho contenido. Para pagar publicidad me dijeron: ``mejor creá más contenido y a partir de ahí pagá publicidad''. Entonces, ahora que me estoy moviendo más en las redes, mi idea es pagar publicidad para que se mueva un poco más.

\noindent\textbf{Tomás Gonçalves:} ¿Tenés más o menos un lugar específico por donde tenías pensado encarar la publicidad?

\noindent\textbf{Morena Cejas:} Sí, más que nada Instagram. En mi opinión, es como que por ahí lo ve la gente.

\noindent\textbf{Micaela Palomino:} Principalmente Instagram, ¿y después tenés TikTok o alguna otra red?

\noindent\textbf{Morena Cejas:} Tengo TikTok, lo uso, pero lo que más subo y lo que más muevo es Instagram. Más allá, también mi círculo, el boca en boca. Por ejemplo, mi cuñado tiene una barbería, entonces lo que hice fue poner un perchero ahí, como para que la gente lo vea. Aprovecho el público, porque mi ropa es unisex, y van desde adolescentes hasta gente grande, y que se mueva ahí.

\noindent\textit{--- Política de precios y promociones ---}

\noindent\textbf{Tomás Gonçalves:} ¿Ajustás precios seguido?

\noindent\textbf{Morena Cejas:} No, la verdad que los precios los mantuve bastante. Son los mismos precios del principio. Ahora hago ofertas, como para que se mueva lo que quedó medio clavado. Mi idea era hacer ropa con buena calidad y más o menos accesible.

\noindent\textbf{Micaela Palomino:} ¿Vos descontás más que nada cuando algo se va a pasar de temporada? ¿O por ejemplo por talles, si ves que hay talles que se mueven más que otros?

\noindent\textbf{Morena Cejas:} Lo hice más que nada por temporadas. Ponele, ahora hace poco fue el Hot Sale y aproveché esos momentos donde todas las marcas hacen descuentos, porque la gente está en el momento de decir: ``bueno, me lo tengo que comprar''.

\noindent\textit{--- Estacionalidad y patrones de venta ---}

\noindent\textbf{Tomás Gonçalves:} Me imagino que notás más ventas en fechas específicas o por temporadas.

\noindent\textbf{Micaela Palomino:} ¿Sentís que hay días o semanas en las que vendés notablemente más?

\noindent\textbf{Morena Cejas:} Sí, en esos momentos como el Hot Sale o el Black Friday. A mí me pasa también que me encanta comprarme cosas; es como que de repente todo va a estar súper barato (que también es un engaño, porque al final es el mismo precio pero parece mucho más barato). Pero sí, en esos momentos yo creo que es donde más se vende, donde más se mueve.

\noindent\textbf{Micaela Palomino:} Y sacando las efemérides como Navidad y Día del Padre, y estos días particulares como Black Friday, después en la diaria, ¿sentís que hay algún día en el que vendés mucho más, o que las primeras semanas del mes se vende más?

\noindent\textbf{Morena Cejas:} Sí, puede ser la primera semana, cuando la gente cobra y aprovecha. Por ahí le tiene ganas a algo y lo compra: un buzo o lo que sea.

\noindent\textit{--- Liquidación de temporada y estrategias de combo ---}

\noindent\textbf{Micaela Palomino:} Entiendo que arrancaste justo en verano. ¿A veces se te ocurre, por ejemplo, que una bermuda se está por pasar de temporada y le ponés promoción, o es más por instinto del momento?

\noindent\textbf{Morena Cejas:} No, sí, lo he hecho. Ahora justamente lo que más estoy liquidando son las primeras remeras. Las remeras se venden siempre porque es algo que usás siempre, pero como fueron de la mano con las bermudas, las liquido juntas. Ahora que hace frío las vendo a mitad de precio para que se las lleven, o hago un combo de remera y bermuda para que se vayan las dos juntas.

\subsubsection*{Bloque 4 · Puntos de dolor y visibilidad del negocio}

\noindent\textit{--- Intuición y falta de datos para decidir ---}

\noindent\textbf{Micaela Palomino:} Ahí veo que el emprendimiento es más análisis tuyo: intuición, ver cuánto te queda, ver cómo lo reponés y qué modelos salieron más.

\noindent\textbf{Morena Cejas:} Claro, literal. Más que nada lo manejo yo porque el emprendimiento es mío. Le pido ayuda a mi papá, que es contador, para que me ayude con los números, las ofertas, cuándo bajar o subir precios, hacer alguna oferta. Pero en general lo manejo yo.

\noindent\textbf{Tomás Gonçalves:} Me imagino que hay cosas de las que te enterás tarde.

\noindent\textbf{Morena Cejas:} Sí, me ha pasado con las remeras. Saqué unas remeras nuevas que estaban mucho más buenas que las primeras. Entonces saqué las segundas y se vendieron más que las primeras.

\noindent\textbf{Micaela Palomino:} ¿Y ahí qué hacés?

\noindent\textbf{Morena Cejas:} Lo que te digo: las primeras remeras las metí con las bermudas, hice un combo de bermuda y remera, esa promo.

\noindent\textit{--- Sorpresas en el comportamiento de venta ---}

\noindent\textbf{Micaela Palomino:} ¿Sentís que tenés visibilidad clara de lo que pasa en tu tienda, o a veces te sorprenden las cosas que suceden?

\noindent\textbf{Morena Cejas:} Sí, me pasó con esas remeras. Cuando empecé la marca dije: hago estos diseños porque me quería quedar en lo seguro. Y al final, las que más me arriesgué vendieron más que las más simples.

\noindent\textbf{Micaela Palomino:} Realmente es por intuición.

\noindent\textbf{Tomás Gonçalves:} Y qué bronca, porque es como ``hago esto, hago lo otro'' a ciegas.

\noindent\textbf{Morena Cejas:} Sí, me pasa ahora: quería hacer un conjunto nuevo para el invierno, aprovechando el clima, y estaba entre dos colores: negro y rojo. Tenía ganas de mandar el rojo, que era más nuevo y distinto. La mitad de la gente me dice: ``hacé el negro por las dudas y después hacé el rojo''. Y muchos me dicen: ``hacé el rojo y arriesgate''. Estoy en ese dilema.

\subsubsection*{Bloque 5 · Validación de la propuesta de valor}

\noindent\textit{--- Presentación de la herramienta y reacción inicial ---}

\noindent\textbf{Micaela Palomino:} Si tuvieras una herramienta que vos entrás y te dice literalmente lo que tenés que hacer ---por ejemplo, analiza las ventas de tu última semana y te dice: ``te conviene reponer tal producto, buzo negro, porque ya sabés que lo vas a vender''---, ¿la pagarías?

\noindent\textbf{Morena Cejas:} Sí. Siendo yo sola la que maneja todo en el emprendimiento, creo que una herramienta así me ayudaría un montón. Como un componente de inteligencia artificial que se fije en las tendencias ---las de mi emprendimiento y las del resto--- y me haga un porcentaje para decirme: ``bueno, ésta hacela de nuevo''. Sí, estaría buenísimo.

\noindent\textbf{Micaela Palomino:} Lo que nosotros habíamos visto era que el punto de dolor es un patrón en todos los emprendimientos: básicamente sos tu propia jefa y aparte tenés que ocupar todos los roles, y a eso le sumás la parte de análisis. Tenés que ver cuándo lanzar una promo, qué hacer con este producto, si pongo cinco o me mando a hacer un rollo entero. Y en parte es por la poca visibilidad que te da la herramienta: te dice ``catorce personas tienen un carrito y no confirmaron la compra'', pero bueno, ¿y qué hago?

\noindent\textbf{Morena Cejas:} Claro, ¿qué producto?

\noindent\textbf{Micaela Palomino:} ¿Cómo sigo? ¿Qué hago con esa información? Nosotros partimos de la premisa de que todas las herramientas dan datos, pero no te dicen qué hacer con eso. Entonces, por intuición tenés que ver qué lográs hacer, y encima que te salga bien.

\noindent\textbf{Morena Cejas:} Claro, concuerdo. Es más, yo pagaba el plan ese de Tienda Nube, y después como tampoco me rindió tanto, dije: esa plata la pongo en publicidad. Listo, no lo pago más.

\noindent\textit{--- Opinión sobre Tienda Nube como plataforma ---}

\noindent\textbf{Morena Cejas:} Yo elegí Tienda Nube porque muchas tiendas la usan, era gratuita y era muy intuitiva para usar. Fue tipo: ``bueno, esta es mi opción''. Pero en el aspecto de análisis era como: ``bueno, ahora me voy a armar un Excel aparte para manejar todo el resto''.

\noindent\textit{--- Disposición a pagar ---}

\noindent\textbf{Tomás Gonçalves:} ¿Cuánto pagabas por el plan de Tienda Nube?

\noindent\textbf{Morena Cejas:} Después del gratuito, el siguiente. Algo de veinte mil pesos, algo así.

\noindent\textbf{Tomás Gonçalves:} Y si viniera con una herramienta de este estilo, ¿más o menos qué rango te sentirías cómoda pagando?

\noindent\textbf{Morena Cejas:} No tengo idea exacta de cuánto pagaría. Depende: si me sirve en ventas, y esta herramienta realmente me ayuda a vender, la pago. Yo no me doy tanta maña con el análisis de stock; mi papá me da una mano porque sabe de números y Excel. Pero una herramienta así me serviría mucho, porque no tengo mucha idea de eso por mi cuenta. Más para decir: ``bueno, esta prenda sirvió, listo, la hago de nuevo''.

\noindent\textbf{Micaela Palomino:} Ya sé que no te la querés jugar con el precio, pero ¿nos decís al menos un máximo de cuánto pagarías?

\noindent\textbf{Morena Cejas:} Tienda Nube, solo por el método de pago y el método de envío, son veinte mil pesos. Entonces, cincuenta mil o sesenta mil pesos por algo que me organice todo y me diga qué hacer, me parece razonable.

\noindent\textit{--- Plataforma integrada versus aplicación separada ---}

\noindent\textbf{Micaela Palomino:} ¿Te gustaría que sea una herramienta aparte o que viva dentro de Tienda Nube?

\noindent\textbf{Morena Cejas:} Puede ser aparte, no tiene por qué ser parte de Tienda Nube. Nos sirve tener distintos puntos de vista.

\noindent\textbf{Micaela Palomino:} La persona que entrevistamos la vez anterior nos contaba que prefería que esté todo en Tienda Nube, le parecía más cómodo tener todo en una sola aplicación.

\noindent\textbf{Morena Cejas:} También puede ser conectar mi cuenta de Tienda Nube con la otra aplicación y que se cargue todo automáticamente. Eso también funcionaría.

\subsubsection*{Bloque 6 · Cierre}

\noindent\textbf{Tomás Gonçalves:} ¿Sentís que hay algo que no te preguntamos?

\noindent\textbf{Morena Cejas:} No sé bien qué necesitaban saber, pero supongo que eso: cómo se maneja Tienda Nube y la poca información que da. Me parece genial que hagan una herramienta para ayudar a los que nos encargamos de todo, básicamente.

\noindent\textbf{Micaela Palomino:} Esta entrevista era para entender cómo se manejan los emprendedores. Y siento que todos se manejan igual: vos tenés un Excel, a mi amiga el novio le hizo una aplicación. Siempre lo terminás tratando fuera de Tienda Nube.

\noindent\textbf{Morena Cejas:} Sí, totalmente. Además, con el Excel tengo que acordarme cuál vendí por promoción y por qué, y se me mezcla también.

\noindent\textbf{Micaela Palomino:} ¿Y en el Excel qué tenés anotado? ¿Tenés el talle, el color?

\noindent\textbf{Morena Cejas:} Tengo todo: la prenda, los talles que hice, la cantidad de cada talle y por color también, el precio de costo de la prenda y el precio al que la vendí o la vendo. Y después cuál vendí en cada modalidad y con qué método de pago.

\noindent\textbf{Micaela Palomino:} Muchas gracias por estar un sábado a las cuatro de la tarde acá con nosotros. Nos sirvió mucho que nos cuentes cómo gestionan los emprendedores las ventas, el stock y los análisis.

\noindent\textbf{Morena Cejas:} Cualquier cosa me dicen, yo no tengo problema.

\begin{center}\textit{--- Fin de la entrevista ---}\end{center}
